
\section[Un caso particular: ser cohomológicamente propio]{Un caso particular de la
propiedad de ser cohomológicamente propio: divisores de cruzamientos normales relativos}
\label{XIII.2}

\setcounter{subsection}{-1}

\subsection{}
\label{XIII.2.0}
Sean $R$ un anillo de valoración discreta de cuerpo de fracciones $K$ y $L$
una $K$-álgebra étale; entonces $L$ es producto directo de un número finito
de cuerpos $L_i$, donde $L_i$ es una extensión étale de~$K$. Si $L'_i$
designa la extensión galoisiana generada por $L_i$ dentro de una clausura
algebraica de~$L_i$, se dice que $L$ está \emph{moderadamente ramificada}
\index{moderadamente ramificada (álgebra)|hyperpage}%
sobre~$R$ si los $L'_i$ son extensiones moderadamente ramificadas en el
sentido de~\Ref{X.3}, \ie si un grupo de inercia~$I_i$ de~$L'_i|K$ tiene
orden primo con la característica residual $p$ de~$R$.

Se sabe que $I_i$ es, en todo caso, extensión de un grupo cíclico de orden
primo con~$p$ por un $p$-grupo. (Esto resulta de
\cite[ch\ptbl IV prop\ptbl 7 cor\ptbl 4]{XIII.5} cuando la extensión residual
de~$R$ es separable. La demostración dada en \loccit se extiende al caso
general de la manera siguiente. Retomemos las hipótesis y la notación de
\loccit, pero sin suponer separable la extensión residual. Sea $H_i$ el
subgrupo del grupo de inercia $G_0$ formado por los elementos $s$ de~$G_0$
tales que $s\pi/\pi\in U^i$ para todo uniformizante $\pi$ de~$A_L$. Se
comprueba entonces que $G_0/H_1$ es un grupo de orden primo con~$p$ y que, para
$i\geqslant 1$, los
\ifthenelse{\boolean{orig}}
{$H^i/H^{i+1}$}
{$H_i/H_{i+1}$}
son $p$-grupos, de donde se deduce el resultado anunciado.)

\subsubsection{}
\label{XIII.2.0.1}
Cuando $R$ es estrictamente local se tiene la caracterización sencilla
siguiente: la $K$-álgebra $L$ está moderadamente ramificada sobre~$R$ si y
solo si los $[L_i:K]$ son primos con~$p$. Además, si $L$ está moderadamente
ramificada sobre~$R$, los $L_i$ son extensiones cíclicas de~$K$. En efecto,
cuando $R$ es estrictamente local, $I_i$ es igual al grupo de Galois
\marginpar{370}
de~$L'_i$ sobre~$K$. Como acaba de recordarse, $I_i$ es extensión de un grupo
cíclico de orden primo con~$p$ por un $p$-grupo. Si se supone que $L'_i$ está
moderadamente ramificada sobre~$R$, entonces $I_i$ es un grupo cíclico de
orden primo con~$p$. Se sigue que $[L_i:K]$ es primo con~$p$ y que $L_i=L'_i$.
Recíprocamente, si $[L_i:K]$ es primo con~$p$, $I_i$ no puede contener ningún
$p$-subgrupo normal no trivial; por tanto, $I_i$ es un grupo cíclico de orden
primo con~$p$, lo que demuestra que $L_i$ está moderadamente ramificada
sobre~$R$.

\subsubsection{}
\label{XIII.2.0.2}
Sean $R$ un anillo de valoración discreta de cuerpo de fracciones $K$, $L$
una $K$-álgebra étale, $\overline{R}$ una localización estricta de~$R$,
$\overline{K}$ su cuerpo de fracciones y
$\overline{L}=L\otimes_K\overline{K}$. Entonces $L$ está moderadamente
ramificada sobre~$R$ si y solo si $\overline{L}$ está moderadamente
ramificada sobre~$\overline{R}$. En efecto, se reduce al caso en que $L$ es
un cuerpo. Escribamos $\overline{L}=\prod_i\overline{L}_i$, donde los
$\overline{L}_i$ son extensiones de~$\overline{K}$; si $L'$ es la
extensión galoisiana generada por $L$ y si
\ifthenelse{\boolean{orig}}
{$L'=L'\otimes_K\overline{K}$,}
{$\overline{L'}=L'\otimes_K\overline{K}$,}
se tiene asimismo una descomposición de~$\overline{L'}$ en producto de
cuerpos,
\ifthenelse{\boolean{orig}}
{$L'=\prod_j\overline{L'_j}$}
{$\overline{L'}=\prod_j\overline{L'_j}$}
y cada $\overline{L}_i$ es una subextensión de al menos uno de los
$\overline{L'_j}$. Como $L'$ es una extensión galoisiana de~$K$, los
$\overline{L'_j}$ son extensiones galoisianas de~$\overline{K}$. Supongamos
que $L$ está moderadamente ramificada sobre~$R$; como el grupo de Galois de
\ifthenelse{\boolean{orig}}
{$\overline{L'_j}|K$}
{$\overline{L'_j}|\overline{K}$}
es isomorfo al grupo de inercia de~$L'|K$, los
\ifthenelse{\boolean{orig}}
{$L'_j$}
{$\overline{L'_j}$}
están también moderadamente ramificados sobre~$\overline{R}$ y, por tanto,
lo mismo sucede con los $\overline{L}_i$. Recíprocamente, supongamos que
$\overline{L}$ está moderadamente ramificada sobre~$\overline{R}$. Para cada
$j$, sea $v_j$ la valoración discreta de~$\overline{L'_j}$ que prolonga la
valoración de~$\overline{K}$, y denotemos también por $v_j$ la valoración
inducida sobre~$L'$. Al variar~$j$, $v_j$ recorre el conjunto de las
valoraciones de
\ifthenelse{\boolean{orig}}
{$L$}
{$L'$}
que prolongan la valoración de~$K$. Sean $G=\Gal(L'|K)$,
$H=\Gal(L'|L)$, $I_j$ el grupo de inercia de~$L'|K$ en $v_j$ y $J_j$ el
grupo de inercia de~$L'|L$ en $v_j$. El grupo
\ifthenelse{\boolean{orig}}
{$I_i$}
{$I_j$}
es extensión de un grupo cíclico de orden primo con~$p$ por un $p$-grupo
$P_j$. Como los $\overline{L}_i$ están moderadamente ramificados
sobre~$\overline{R}$, $I_j/J_j$ tiene orden primo con~$p$ y, por tanto,
$P_j\subset J_j$. En consecuencia, el grupo $H$ contiene todos los $P_j$ y,
por tanto, también el grupo $P$ engendrado por
\marginpar{371}
los $P_j$ al variar~$j$. Pero el grupo~$P$ es invariante en~$G$, pues un
automorfismo interior de~$G$ permuta los~$I_j$ entre sí y, por tanto, también
los $P_j$ entre sí. Se sigue que $P$ es un subgrupo de~$H$ normal en~$G$ y,
por consiguiente, puesto que $L'$ es la extensión galoisiana generada por
$L$, que $P=1$, lo que demuestra que $L$ está moderadamente ramificada
sobre~$R$.

Más generalmente, sea $R\to R'$ un homomorfismo de anillos de valoración
discreta tal que la imagen de un uniformizante $\pi$ de~$R$ sea un
uniformizante $\pi'$ de~$R'$ y que la extensión residual $\kres(R')$ sea una
extensión separable de~$\kres(R)$. Sean $K$ el cuerpo de fracciones de~$R$,
$K'$ el cuerpo de fracciones de~$R'$, $L$ una $K$-álgebra étale y
$L'=L\otimes_KK'$. Entonces $L$ está moderadamente ramificada sobre~$R$ si y
solo si $L'$ está moderadamente ramificada sobre~$R'$. En efecto, puede
suponerse que $R$ y $R'$ son estrictamente locales. Según~\Ref{XIII.2.0.1},
basta probar que, cuando $L$ es un cuerpo, también lo es~$L'$. Sean
$\tilde{R}$ la normalización de~$R$ en~$L$, $\tilde{\pi}$ un
uniformizante de~$\tilde{R}$ y
$\tilde{R}'=\tilde{R}\otimes_RR'$. Como la extensión
$\kres(\tilde{R})|\kres(R)$ es radicial y la extensión
$\kres(R')|\kres(R)$ es étale,
$\kres(\tilde{R})\otimes_{\kres(R)}\kres(R')$ es un cuerpo [EGA IV 4.3.2
y 4.3.5]. Esto demuestra que
\ifthenelse{\boolean{orig}}
{$R'$}
{$\tilde{R}'$}
es un anillo local y, como $\pi$ tiene por imagen $\pi'$ en~$R'$, se tiene
\ifthenelse{\boolean{orig}}
{$\kres(\tilde{R}')=R'/(\tilde{\pi})$;}
{$\kres(\tilde{R}')=\tilde{R}'/(\tilde{\pi})$;}
por consiguiente,
\ifthenelse{\boolean{orig}}
{$R'$}
{$\tilde{R}'$}
es un anillo de valoración discreta
\cite[ch\ptbl I \S 2 prop\ptbl 2]{XIII.5} y, por tanto, $L'$ es un cuerpo.

\subsubsection{}
\label{XIII.2.0.3}
Reduciendo al caso estrictamente local, se ve que una subálgebra de un
álgebra moderadamente ramificada está moderadamente ramificada; que el
producto tensorial de dos álgebras moderadamente ramificadas está
moderadamente ramificado; que un álgebra moderadamente ramificada sigue
siéndolo tras extender el anillo de valoración discreta; y que un álgebra
que pasa a estar moderadamente ramificada tras una extensión moderadamente
ramificada está moderadamente ramificada.
